% !TEX root = ../notes.tex

\documentclass[../notes.tex]{subfiles}

\begin{document}

\section{May 4}
Ok, let's start.

\subsection{Vector Bundles with Flat Connection}
We are interested in vector bundles with a flat connection.
\begin{definition}[flat connection]
	Fix a smooth algebraic variety $X$ over an algebraically closed field $k$. Then a \textit{vector bundle} is a locally free sheaf of finite rank. A \textit{flat connection} on $\mc E$ is a map $\nabla\colon\mc E\to\mc E\otimes\Omega^1$ such that $\nabla\land\nabla=0$.
\end{definition}
\begin{remark}
	One can view vector fields $v$ on $X$ as derivations $\OO_X\to\OO_X$. (One can also work locally on open subsets.) Then $\nabla$ induces an action by vector fields on sections of $\mc E$, denoted $\nabla_v$, and flatness amounts to asking for $\nabla_{[v,w]}=[\nabla_v,\nabla_w]$. Thus, the Lie algebra $\op{Vect}(X)$ acts on sections of $\mc E$.
\end{remark}
\begin{remark}
	The action by the Lie algebra on sections satisfies
	\[\nabla_v(f\sigma)=f\nabla_v(\sigma)+v(f)\sigma.\]
\end{remark}
Here are some examples, from complex analysis.
\begin{example}
	With $k=\CC$, the vector bundle is an analytic bundle. Choosing local trivializations gives rise to a local system, which is the sheaf of locally constant sections which are parallel to $\nabla$ (also called ``flat sections''). The data of such a thing is equivalent to the data of a representation of $\pi_1(X;x_0)$ when $X$ is connected.
\end{example}
\begin{example}
	Given a smooth projective morphism $f\colon Y\to X$ of smooth projective varieties over $\CC$, one has a local system $\mathrm R^if_*^{\mathrm{an}}\CC$ on $X$ whose fiber over $x\in X$ is $\mathrm H^i(Y_x;\CC)$ (by some base-change theorems). Hodge theory endows $\mathrm H^i(Y_x;\CC)$ with a pure Hodge structure
	\[\mathrm H^i(Y_x;\CC)=\bigoplus_{p+q=i}\mathrm H^q(\Omega^p_{Y_x};\CC).\]
	(There is also an integral structure on cohomology, which we will ignore for now.) This decomposition turns out to be induced by vector bundle with flat connection, known as the Gauss--Manin connection. Namely, for a given $n$, one can define $\mathrm R^i_{\le n}f_*\CC$ to have fibers $\bigoplus_{q\le n}\mathrm H^q(\Omega^p_{Y_x};\CC)$, and this is an algebraic sub-bundle. Furthermore, $\nabla$ increases degree in $n$ by one.
\end{example}
For more examples, we can expand our definition.
\begin{remark}
	The exact same definition allows us to define quasicoherent sheaves with flat connection. Using the relation $\nabla_v(f\sigma)=f\nabla_v(\sigma)+v(f)\sigma$, one finds that such objects have the same data as a quasicoherent $\OO_X$-module over the sheaf of rings $\mc D_X$. Here, $\mc D_X$ is the algebra generated by functions $\OO_X$ and vector fields, satisfying the relations $f\cdot v=fv$ and $v\cdot w-w\cdot v=[v,w]$ and $v\cdot f=f\cdot v=v(f)$. To check this, one needs to use the relation expanding $\nabla_v(f\sigma)$.
\end{remark}
\begin{example}
	For $X=\AA^n$, we find that quasicoherent sheaves with flat connections are the same as modules over the Weyl algebra
	\[\mc D_X=\frac{k\langle x_1,\ldots,x_n,\del_1,\ldots,\del_n\rangle}{([\del_i,x_j]-1_{i=j})_{i,j}}.\]
\end{example}

\subsection{Quantiation}
One can also quantize $\mc D$, using the Rees construction.
\begin{remark}
	Note that $\mc D$ admits a natural filtration by the order of the differential form. Thus, $\op{gr}\mc D_X$ is the same as functions on $T^*X$ because now the vector fields commute with functions.
\end{remark}
\begin{notation}
	Fix a smooth projective variety $X$ over an algebraically closed field $k$. Then we define $\mc D_{\hbar}$ to be a sheaf of $k[\hbar]$-modules generated by functions and vector fields on $X$, satisfying the relations $f\cdot v=fv$ and $v\cdot w-w\cdot v=\hbar[v,w]$ and $v\cdot f-f\cdot v=\hbar v(f)$.
\end{notation}
\begin{example}
	By definition, $\mc D_1=\mc D$ and $\mc D_0=\OO_{T^*X}$.
\end{example}
\begin{example}
	A quasicoherent sheaf on $T^*X$ is basically a Higgs field. Then the relevant flat connection $\overline\nabla\colon\mc E\to\mc E\otimes\Omega^1$ is a morphism of $\OO$-modules satisfying $\overline\nabla^2=0$.
\end{example}
\begin{example}
	Fix a smooth projective morphism $f\colon X\to Y$ of complex varieties. Then $\mc E=\mathrm R^*f_*\CC$ as a vector bundle with the flat connection given by the Gauss--Manin connection. Taking the associated graded $\op{gr}\mc E$ gives the Higgs field
	\[\bigoplus_n\mc E_{\le n}/\mc E_{\le n-1}.\]
\end{example}
These sorts of filtrations will be interesting to us.
\begin{remark}
	More generally, if $\mc M$ is a $\mc D$-module equipped with a filtration $\{\mc M_{\le i}\}_i$ by $\OO$-submodules, and for which the vector field action increases filtration degree by at most one, then $\op{gr}\mc M$ is a quasicoherent sheaf on $T^*X$. In fact, one can find a global module $\widetilde{\mc M}$ over $\mc D_{\hbar}$ which specializes to both $\mc M$ and $\op{gr}\mc M$.
\end{remark}
\begin{definition}[good filtration]
	Fix a smooth algebraic variety $X$ over an algebraically closed field $k$. Then a \textit{good filtration} on a $\mc D$-module $\mc M$ is an exhaustive filtration $\{\mc M_{\le i}\}$ for which $\op{gr}\mc M$ is a coherent sheaf on $T^*X$.
\end{definition}
This perspective will allow us to clarify some things. For example, let's assume that there are well-defined notions of pulling back and pushing forward vector bundles with flat connection and more generally $\mc D_{\hbar}$-modules.
\begin{example}
	Let $f\colon Y\to X$ be a smooth morphism. Then there is a diagram
	% https://q.uiver.app/#q=WzAsMyxbMCwwLCJUXiooWFxcdGltZXNfWFkpIl0sWzEsMCwiVF4qWSJdLFswLDEsIlReKlgiXSxbMCwxLCJkZl4qIiwwLHsic3R5bGUiOnsidGFpbCI6eyJuYW1lIjoiaG9vayIsInNpZGUiOiJ0b3AifX19XSxbMCwyLCJcXG9we3ByfV8xIiwyXV0=&macro_url=https%3A%2F%2Fraw.githubusercontent.com%2FdFoiler%2Fnotes%2Fmaster%2Fnir.tex
	\[\begin{tikzcd}[cramped]
		{T^*(X\times_XY)} & {T^*Y} \\
		{T^*X}
		\arrow["{df^*}", hook, from=1-1, to=1-2]
		\arrow["{\op{pr}_1}"', from=1-1, to=2-1]
	\end{tikzcd}\]
	where $df^*$ is given by the dual action. There is a quantization of this diagram to $\mc D$-modules given by a single $(f^\sharp\mc D_X,\mc D_Y)$-bimodule $\mc D_{Y\to X}$, which encodes the construction of pulling back along $df$ and pushing forward along $\op{pr}_1$. Namely, at $\hbar=1$, the functor $\mc M\mapsto\mathrm R\op{pr}_\sharp(D_{X\to Y}\otimes^{\mathrm L}_{\mc D_Y}\mc M)$. At $\hbar=0$, we recover the pull-push functor. This functor provides a spectral sequence, and Hodge theory turns out to imply that
	\[\op{gr}\mathrm R^if_*\CC=\mathrm H^i(\op{pr}_*df^*\OO_Y).\]
\end{example}

\subsection{Positive Characteristic}
We now pass to positive characteristic $p$. The chief new features are downstream of the Frobenius. For example, the functions $f^p$ are killed by any vector field, and $(-)^p$ sends derivations to derivations.
\begin{notation}
	Fix a smooth algebraic variety $X$ over a perfect field of positive characteristic $p$. Then we define $(-)^{[p]}$ to send the vector field $v$ with associated derivation $\delta$ to the vector field with associated derivation $\delta^p$.
\end{notation}
\begin{remark}
	For any vector field $v$, we see that $v^p-v^{[p]}$ acts by zero on functions, but it need not vanish. Instead, it lives in the center of $\mc D_X$.
\end{remark}
\begin{example}
	With $X=\AA^n_k$, one can check that $\del_i^{[p]}=0$, so $\del_i^p$ is central. Also, $x_i^p$ is central, and one can check that the center is generated by these $p$th powers.
\end{example}
This center reintroduces $T^*X$.
\begin{remark}
	In general, it turns out that $Z(\mc D_X)$ is $\OO_{T^*X^{(1)}}$, where $X^{(1)}$ denotes the Frobenius twist $X\times_{k,\mathrm{Frob}}k$. In fact, $\mc D_X$ is an Azumaya algebra over $T^*X^{(1)}$ of rank $p^{2\dim X}$, and its order in the Brauer group is $p$. As a sketch, one can find $\mc D_X$ as the image of the canonical Liouville $1$-form $\theta=\sum_ip_i\,dq_i$ on $T^*X$ along a natural map $\mathrm H^0(X;\Omega^1)\to\op{Br}(T^*X)[p]$.
\end{remark}
\begin{example}
	Fix a field $k$ of positive characteristic $p$, and consider $X=\AA^1$. Then we claim that the Weyl algebra is $k\langle x,\del\rangle/([\del,x]-1)$. Set $\alpha\coloneqq x^p$ and $\beta\coloneqq\delta^p$ for brevity, and then we can take a further quotient to requiring that $\alpha,\beta\in k$. One can reduce to the case where $\alpha=\beta=0$, where one finds that the Weyl algebra is the ring of $p\times p$ matrices acting on the vector space $k[x]/\left(x^p\right)$. This is analogous to the Stone--von Neumann theorem.
\end{example}
\begin{remark}
	It turns out that $\mc D_X|_{X^{(1)}}$ is isomorphic to $\op{End}_{\OO_{X^{(1)}}}(\op{Fr}_*\OO_X)$, and this latter Azumaya algebra splits on the zero section of $T^*X^{(1)}$. Notably, we see that any $\mc D$-module is naturally a module on $T^*X^{(1)}$, and if the Azumaya algebra has scheme-theoretic support on the zero section $X^{(1)}\subseteq T^*X^{(1)}$, then $\mc M$ is isomorphic to its Frobenius twist.
\end{remark}
\begin{remark}
	If $X$ admits a lift to $W_2(k)$, then the Azumaya algebra splits on a $p$-neighborhood of the zero section.
\end{remark}
The action of $Z(\mc D_X)$ is related to the so-called $p$-curvature.
\begin{defihelper}[$p$-curvature] \nirindex{pcurvature@$p$-curvature}
	Fix a smooth algebraic variety $X$ over a field $k$ of positive characteristic $p$. Given a vector bundle $\mc E$ with flat connection $\nabla$, we define the \textit{$p$-curvature} along a vector field $v$ to be the endomorphism $\nabla_v^p-\nabla_{v^{[p]}}$ outputting sections of $\op{Fr}^*\Omega^1\otimes\op{End}\mc E$.
\end{defihelper}
Notably, this notion of curvature only remembers a single vector field: it can already be interested in just one dimension!

We will be interested in how characteristic zero compares to positive characteristic.
\begin{remark}
	If $X$ is defined over a number field $K$, then all but finitely many of its reductions to $\OO_K/\mf p$ will continue to be smooth algebraic varieties. Furthermore, if $(\mc E,\nabla)$ is a vector bundle with flat connection, again all but finitely many of the reductions continue to let this be a vector bundle with flat connection.
\end{remark}
For example, one could try to pass information between the two settings.
\begin{conj}[Grothendieck]
	Fix a smooth algebraic variety $X$ over a number field $K$, and choose a vector bundle $\mc E$ with flat connection $\nabla$ defined over $K$. Suppose that $(\mc E_{k(\mf p)},\nabla_{k(\mf p)})$ has zero $p$-curvature for all but finitely many $\mf p$. Then $(\mc E_\CC,\nabla_\CC)$ has finite monodromy.
\end{conj}
\begin{remark}
	If $(\mc E_\CC,\nabla_\CC)$ has regular singularities, then the converse is true: the local system will trivialize by some finite analytic covering, but such a finite analytic cover must be algebraic. Because $(\mc E_\CC,\nabla_\CC)$ has regular singularities, the Riemann--Hilbert correspondence implies that $(\mc E,\nabla)$ itself trivializes after such a finite covering, which is enough to cause the $p$-curvature to vanish.
\end{remark}
\begin{remark}
	The choice of global model spreading out $X$ can change the set of finite places with nonzero $p$-curvature.
\end{remark}
\begin{remark}
	Much is known about this conjecture. It was shown in rank one in the 1980s by Honda and the Chudnovsky brothers. Katz showed this for Gauss--Manin connections in 1972, which we will discuss next time.
\end{remark}
Alternatively, one can try to understand the $p$-support as $p$ varies. The following is due to many people.
\begin{theorem}
	Fix a smooth algebraic variety $X$ over a number field $K$, and choose a vector bundle $\mc E$ with flat connection $\nabla$ defined over $K$. Then for all but finitely many $\mf p$, the $p$-support of $(\mc E_{k(\mf p)},\nabla_{k(\mf p)})$ is Lagrangian.
\end{theorem}
Here are some examples.
\begin{example}
	With $\mc E=\OO$, one can take $\nabla=d+\omega$, where $\omega$ is some $1$-form. Then the $p$-curvature is $\omega-C(\omega)$, where $C$ is the Cartier isomorphism. For example, taking $X=\AA^1$ and $\omega=dt$, one has $C(\omega)=0$. Similarly, taking $X=\mathbb G_m$ and $\omega=\lambda\,dt/t$, then $C(\omega)=(\lambda^p-\lambda)\,dT/T$, where $T$ is the local coordinate on $X^{(1)}$.
\end{example}

\end{document}